\begin{abstract}
The dip test of Hartigan and Hartigan (1985) is the standard nonparametric test of unimodality. It measures how far the empirical distribution function of a sample lies from the nearest unimodal distribution function, that is, from a distribution function that curves upward (is convex) until the peak of the density and curves downward (is concave) after it. Its critical values are taken from the uniform distribution, on the grounds that among unimodal laws the uniform produces the largest dip. We show that this is false. Consider a density that is flat over most of its range and then jumps up to a narrow peak at one end. Because the peak sits at the end of the flat stretch, every good unimodal fit to the sample must keep curving upward over the whole flat stretch, and it can no longer bend to follow random bumps in the data. In large samples the dip then behaves like half the largest gap between a Brownian bridge and its concave majorant, a quantity whose law was found by Balabdaoui and Pitman (2011), and this quantity is stochastically larger than the uniform limit of the dip. With 90 percent of the mass on the flat stretch, the nominal 5 percent test rejects about one sample in four, however large the sample.

We prove that over atomless unimodal laws with bounded support and a finite density at the mode, the worst limiting rejection rate of the standard 5 percent test lies between 28.4 and 33.2 percent, and between 12.7 and 16.8 percent at the 1 percent level; the width of these intervals comes from a certified numerical bracket for the limiting uniform critical value. Rejecting when the dip exceeds $c_\alpha/\sqrt n$, where $c_\alpha$ is half the $(1-\alpha)$ quantile of the Balabdaoui--Pitman gap ($c_{0.05}\approx0.7314$ and $c_{0.01}\approx0.8599$), restores the nominal level on that class in the limit, and no deterministic sequence of cutoffs with a smaller limit does. A computer-assisted argument, based on Monte Carlo certificates with a stated error probability, shows that the same intervals hold when the worst law is chosen separately at each sample size among all continuous unimodal laws. We also give certified cutoffs for small and moderate samples, and we measure the price of the correction: on symmetric normal mixtures the corrected test needs about twice as many observations as the uniform-calibrated one to reach the same power.
\end{abstract}
